\chapter{Convex Sets}

A set is convex if it contains the entire straight line
segment between any two of its points: no holes, no dents, no detached
pieces. In economics this single geometric idea describes budget sets,
production possibility sets, and the feasible regions of linear programs.
Its payoff for optimization is structural: on a convex feasible set a convex
problem has no spurious local optima: every local optimum is automatically
global, and first-order conditions become sufficient rather than merely
necessary.

This chapter develops the geometry and algebra of convex sets: the
line-segment criterion, weighted averages (convex combinations), the standard
convex sets of the plane, the operations that preserve convexity (and the
one that famously does not) and the contour sets of functions, which build
the bridge to concavity and quasiconcavity in the chapters that follow.

\section{Core Rules}

\begin{rulebox}{Definition (line-segment criterion)}
A set $S\subseteq\R^{n}$ is \emph{convex} if for all $x,y\in S$ and all
$\lambda\in[0,1]$,
\[
  \lambda x+(1-\lambda)y \;\in\; S .
\]
As $\lambda$ slides from $0$ to $1$, the point $\lambda x+(1-\lambda)y$
traces the segment from $y$ to $x$; convexity says the whole segment stays
inside $S$. The empty set, single points, lines, and segments are all convex.
\end{rulebox}

\begin{rulebox}{Convex combinations and the convex hull}
A \emph{convex combination} of points
$x_{1},\dots,x_{k}\in\R^{n}$ is any weighted average
\[
  \sum_{i=1}^{k}\lambda_{i}x_{i},
  \qquad \lambda_{i}\ge 0,\quad \sum_{i=1}^{k}\lambda_{i}=1 .
\]
The case $k=2$, $\lambda_{1}=\lambda_{2}=\tfrac12$ is the midpoint. A set is
convex if and only if it is closed under \emph{all} finite convex
combinations of its points. The \emph{convex hull}
$\operatorname{conv}(A)$ of a set $A$ is the set of all finite convex
combinations of points of $A$, equivalently, the smallest convex set
containing $A$. The hull of finitely many points is a polygon whose vertices
are among those points.
\end{rulebox}

\begin{rulebox}{Standard convex sets in the plane}
\begin{itemize}[leftmargin=1.6em,itemsep=0.3ex]
  \item \textbf{Half-planes} $\{(x,y): ax+by\le c\}$ (also with $\ge$, $<$,
        $>$, or $=$). Linear expressions pass through weighted averages:
        \[
          a\bigl(\lambda x_{1}+(1-\lambda)x_{2}\bigr)
          +b\bigl(\lambda y_{1}+(1-\lambda)y_{2}\bigr)
          =\lambda(ax_{1}+by_{1})+(1-\lambda)(ax_{2}+by_{2})
          \le \lambda c+(1-\lambda)c=c .
        \]
  \item \textbf{Budget sets}
        $\{(x_{1},x_{2}): p_{1}x_{1}+p_{2}x_{2}\le m,\ x_{1},x_{2}\ge 0\}$
        with $p_{1},p_{2},m>0$: the intersection of three half-planes, hence
        convex, a triangle with corners $(0,0)$, $(m/p_{1},0)$,
        $(0,m/p_{2})$.
  \item \textbf{Closed disks} $\{(x,y): (x-x_{0})^{2}+(y-y_{0})^{2}\le
        r^{2}\}$: convex by the triangle inequality for the norm.
  \item \textbf{Polyhedra}: intersections of finitely many half-planes
        (triangles, quadrilaterals, feasible regions of linear programs).
        Convex because intersections preserve convexity.
\end{itemize}
\end{rulebox}

\begin{gallery}
\begin{tikzpicture}
  \begin{axis}[geoaxis, xmin=-2.6, xmax=2.6, ymin=-2.6, ymax=2.6,
      xtick={-2,2}, ytick={-2,2}, title={convex: a disc}]
    \addplot[draw=none, name path=up, domain=-2:2, samples=201]{sqrt(4-x^2)};
    \addplot[draw=none, name path=lo, domain=-2:2, samples=201]{-sqrt(4-x^2)};
    \region{up}{lo}{}
    \addplot[domain=0:360, samples=181]({2*cos(x)},{2*sin(x)});
    \draw[black, dashed, thick] (axis cs:-1.5,0.8) -- (axis cs:1.5,-0.8);
    \addplot[only marks, mark=*, mark size=1.1pt, black, forget plot]
      coordinates {(-1.5,0.8) (1.5,-0.8)};
  \end{axis}
\end{tikzpicture}\hfill
\begin{tikzpicture}
  \begin{axis}[geoaxis, xmin=-2.6, xmax=2.6, ymin=-2.6, ymax=2.6,
      xtick={-2,2}, ytick={-2,2}, title={not convex: an annulus}]
    \addplot[draw=none, name path=o1, domain=-2:2, samples=201]{sqrt(4-x^2)};
    \addplot[draw=none, name path=o2, domain=-2:2, samples=201]{-sqrt(4-x^2)};
    \region{o1}{o2}{}
    \addplot[draw=none, name path=i1, domain=-1:1, samples=201]{sqrt(1-x^2)};
    \addplot[draw=none, name path=i2, domain=-1:1, samples=201]{-sqrt(1-x^2)};
    \addplot[white, forget plot, sharp plot] fill between[of=i1 and i2];
    \addplot[domain=0:360, samples=181]({2*cos(x)},{2*sin(x)});
    \addplot[domain=0:360, samples=181]({cos(x)},{sin(x)});
    \draw[black, dashed, thick] (axis cs:-1.5,0) -- (axis cs:1.5,0);
    \addplot[only marks, mark=*, mark size=1.1pt, black, forget plot]
      coordinates {(-1.5,0) (1.5,0)};
  \end{axis}
\end{tikzpicture}\hfill
\begin{tikzpicture}
  \begin{axis}[geoaxis, xmin=-2.6, xmax=2.6, ymin=-0.6, ymax=4.6,
      xtick={-2,2}, ytick={2,4}, title={not convex: $y\le x^{2}$}]
    \addplot[draw=none, name path=par, domain=-2.1:2.1, samples=201]{x^2};
    \addplot[draw=none, name path=bot, domain=-2.1:2.1, samples=2]{-0.55};
    \region{par}{bot}{}
    \addplot[domain=-2.15:2.15]{x^2};
    \draw[black, dashed, thick] (axis cs:-1.5,2.25) -- (axis cs:1.5,2.25);
    \addplot[only marks, mark=*, mark size=1.1pt, black, forget plot]
      coordinates {(-1.5,2.25) (1.5,2.25)};
  \end{axis}
\end{tikzpicture}
\gcap{The line-segment test. In the first set every chord stays inside; in
the other two a chord between two points of the set leaves it, so the sets
are not convex.}
\end{gallery}

\begin{rulebox}{Operations that preserve convexity: and one that does not}
Let $S$, $S_{1}$, $S_{2}$ be convex. Then the following are convex:
\begin{itemize}[leftmargin=1.6em,itemsep=0.3ex]
  \item \textbf{Intersections} $\bigcap_{i}S_{i}$, of \emph{any} (even
        infinite) family of convex sets;
  \item \textbf{scalar multiples} $\alpha S=\{\alpha x : x\in S\}$ for
        $\alpha\in\R$;
  \item \textbf{Minkowski sums} $S_{1}+S_{2}=\{x+y : x\in S_{1},\ y\in
        S_{2}\}$;
  \item \textbf{affine images} $\{Ax+b : x\in S\}$.
\end{itemize}
\textbf{Warning:} the \emph{union} of convex sets is in general \emph{not}
convex. Example: $\{x\le -1\}\cup\{x\ge 1\}$ contains $-1$ and $1$ but not
their midpoint $0$. Complements and set differences of convex sets are
generally not convex either.
\end{rulebox}

\begin{rulebox}{Contour sets: the bridge to quasiconcavity}
For $f:\R^{n}\to\R$ and $c\in\R$ define the \emph{upper contour set}
$U_{c}=\{x: f(x)\ge c\}$ and the \emph{lower contour set}
$L_{c}=\{x: f(x)\le c\}$. The same two sets are also written
$B(c)=U_{c}$ \de{(Superniveaumenge, obere Konturmenge)} and
$W(c)=L_{c}$ \de{(Subniveaumenge, untere Konturmenge)}, and indexed by a
\emph{point} rather than a level: $B(P)=B(f(P))$, $W(P)=W(f(P))$. Both
notations are in common use; \chref{derivatives}{9} introduces them with
the level curves.
\begin{itemize}[leftmargin=1.6em,itemsep=0.3ex]
  \item If $f$ is \textbf{convex}, every lower contour set $L_{c}$ is
        convex.
  \item If $f$ is \textbf{concave}, every upper contour set $U_{c}$ is
        convex.
  \item $f$ is called \textbf{quasiconcave} precisely when \emph{every}
        upper contour set is convex (quasiconvex: every lower contour set
        convex). This is \emph{weaker} than concavity: $u(x_{1},x_{2})
        =x_{1}x_{2}$ on the positive quadrant is quasiconcave but not
        concave.
\end{itemize}
In consumer theory, ``indifference curves convex to the origin'' is exactly
the statement that the upper contour sets of utility are convex.
\end{rulebox}

\begin{rulebox}{Verifying convexity algebraically}
To \emph{prove} that $S=\{x: g(x)\le c\}$ is convex: take arbitrary $x,y\in
S$ and $\lambda\in[0,1]$, and show $g(\lambda x+(1-\lambda)y)\le c$ using
only $g(x)\le c$ and $g(y)\le c$. For linear $g$ this is automatic (see the
half-plane computation above). Two workhorses:
\begin{itemize}[leftmargin=1.6em,itemsep=0.3ex]
  \item the \textbf{square inequality}
        \[
          \bigl(\lambda a+(1-\lambda)b\bigr)^{2}
          \;\le\;\lambda a^{2}+(1-\lambda)b^{2},
        \]
        valid for all $a,b\in\R$, $\lambda\in[0,1]$, since the difference of
        the two sides is $\lambda(1-\lambda)(a-b)^{2}\ge 0$;
  \item the \textbf{graph shortcut}: $\{(x,y): y\ge g(x)\}$ is convex when
        $g''\ge 0$ ($g$ convex), and $\{(x,y): y\le g(x)\}$ is convex when
        $g''\le 0$ ($g$ concave).
\end{itemize}
To \emph{disprove} convexity: exhibit two points of $S$ whose midpoint
violates the defining condition, a single counterexample suffices.
\end{rulebox}

\begin{rulebox}{Strictly convex sets}
A convex set $S\subseteq\R^{n}$ is \emph{strictly convex}
\de{streng konvex} if for all $x\neq y$ in $S$ and all $\lambda\in(0,1)$
the point $\lambda x+(1-\lambda)y$ lies in the \emph{interior} of $S$:
every chord, apart from its two endpoints, runs strictly inside the set.
For a closed convex set this says that the boundary contains no straight
piece: no flat edges, no flat faces.
\begin{itemize}[leftmargin=1.6em,itemsep=0.3ex]
  \item \textbf{Strictly convex:} discs and balls, filled ellipses, the
        region $\{y\ge g(x)\}$ above a strictly convex curve, the region
        $\{xy\ge c,\ x>0\}$ beyond a hyperbola branch ($c>0$).
  \item \textbf{Convex, not strictly:} every polygon and polyhedron
        (triangle, square, cube), half-planes, a half-disc (its diameter is
        flat), a cylinder or a cone (straight lines on the surface). Test:
        two points on the same flat piece of boundary have their midpoint on
        the boundary, not in the interior.
  \item \textbf{Not convex at all:} a notched or dented body, an annulus,
        a torus.
\end{itemize}
In consumer theory an upper contour set $\{u\ge c\}$ is strictly convex
exactly when the indifference curve $u=c$ has no straight segment; this is
what ``strictly convex preferences'' means, and it is the geometric reason
why the utility-maximising bundle on a budget line is unique.
\end{rulebox}

\section{Worked Example}

\begin{worked}{Which of these are convex? The budget set, the circle $x^{2}+y^{2}=1$, the region $\{xy\ge1,\ x>0\}$, and a union of two discs}
\wstep{Step 1. State the criterion you will actually use.}
$S$ is convex if for all $\mathbf a,\mathbf b\in S$ and all
$\lambda\in[0,1]$,
\[
  \lambda\mathbf a+(1-\lambda)\mathbf b\in S .
\]
In words: the segment joining any two points of $S$ stays inside $S$. To
prove convexity you must handle \emph{all} pairs; to disprove it you need
only \emph{one} pair. That asymmetry decides which way to attack.

\wstep{Step 2. The budget set --- convex, and for a reusable reason.}
\[
  B=\{(x,y):p_{x}x+p_{y}y\le M,\ x\ge0,\ y\ge0\}
\]
is the intersection of three half-planes. Each half-plane
$\{\mathbf x:\mathbf a\cdot\mathbf x\le c\}$ is convex, because
\[
  \mathbf a\cdot\bigl(\lambda\mathbf x_{1}+(1-\lambda)\mathbf x_{2}\bigr)
  =\lambda(\mathbf a\cdot\mathbf x_{1})+(1-\lambda)(\mathbf a\cdot\mathbf x_{2})
  \le\lambda c+(1-\lambda)c=c ,
\]
and \textbf{an intersection of convex sets is convex}. So $B$ is convex
without any picture. Every feasible set defined by linear inequalities is
convex for exactly this reason --- which is why linear programming in
\volref{linprog} is as well behaved as it is.

\wstep{Step 3. The circle --- not convex, and the reason is instructive.}
$C=\{(x,y):x^{2}+y^{2}=1\}$ is the \emph{curve}, not the disc. Take
$\mathbf a=(1,0)$ and $\mathbf b=(-1,0)$, both in $C$. Their midpoint is
$(0,0)$, and $0^{2}+0^{2}=0\neq1$. So $C$ is \textbf{not convex}.
\wtrap{A set defined by an \emph{equality} is convex only if the equality is
linear. $\{x+y=1\}$ is a line and is convex; $\{x^{2}+y^{2}=1\}$ is a circle
and is not. Read the relation symbol before anything else: $\le$ and $\ge$
give solid regions, $=$ gives a thin boundary that is almost never convex.
The closed \emph{disc} $x^{2}+y^{2}\le1$ is, of course, convex.}

\wstep{Step 4. $\{xy\ge1,\ x>0\}$ --- convex, proved algebraically.}
Take $(x_{1},y_{1})$ and $(x_{2},y_{2})$ with all coordinates positive and
$x_{i}y_{i}\ge1$. Write $\mu=1-\lambda$ and expand the product at the
intermediate point:
\[
  (\lambda x_{1}+\mu x_{2})(\lambda y_{1}+\mu y_{2})
  =\lambda^{2}x_{1}y_{1}+\mu^{2}x_{2}y_{2}+\lambda\mu\bigl(x_{1}y_{2}+x_{2}y_{1}\bigr).
\]
By the arithmetic--geometric mean inequality,
\[
  x_{1}y_{2}+x_{2}y_{1}\ge2\sqrt{x_{1}y_{2}x_{2}y_{1}}
  =2\sqrt{(x_{1}y_{1})(x_{2}y_{2})}\ge2 .
\]
Hence the product is at least
$\lambda^{2}+\mu^{2}+2\lambda\mu=(\lambda+\mu)^{2}=1$, and the intermediate
point is in the set. \textbf{Convex.} \checkmark

Geometrically this is the region above a hyperbola --- and a hyperbola in
the positive quadrant is a convex curve, so ``above'' is the right side.

\wstep{Step 5. A union of two discs --- not convex, even though each piece is.}
Take the discs of radius $1$ centred at $(0,0)$ and at $(5,0)$. The points
$(0,0)$ and $(5,0)$ both belong to the union, but their midpoint
$(2.5,0)$ is in neither disc.
\[
  \boxed{\text{Intersections of convex sets are convex; unions are not.}}
\]
\wtrap{This asymmetry matters in practice. A constraint written with
``and'' ($g_{1}\le0$ \emph{and} $g_{2}\le0$) preserves convexity; one
written with ``or'' destroys it, and with it every sufficiency theorem in
this volume. Integer restrictions and fixed-cost ``either produce nothing or
at least $q_{0}$'' conditions are exactly of the ``or'' kind, which is why
they need entirely different methods.}

\wstep{Step 6. The operations that are safe.}
\begin{center}\footnotesize
\begin{tabular}{ll}
  \hline
  operation & preserves convexity?\\
  \hline
  intersection (any number) & yes\\
  scaling $cS$, translation $S+\mathbf v$ & yes\\
  sum $S_{1}+S_{2}=\{\mathbf a+\mathbf b\}$ & yes\\
  image under a linear map & yes\\
  union & \textbf{no}\\
  complement & \textbf{no}\\
  \hline
\end{tabular}
\end{center}

\wstep{Step 7. Where this is heading.}
Convexity of the \emph{feasible set} is one of the two hypotheses that turn
a first-order condition into a proof of a global optimum; concavity of the
\emph{objective} is the other. \chref{optimization}{3} supplies the second, and
\chref{optimization}{10} shows what fails when either is missing. Contour sets provide the
bridge: the superlevel sets $\{U\ge\bar u\}$ of a utility function are
convex precisely when $U$ is quasiconcave, and that --- not concavity --- is
what economics actually assumes about preferences.
\end{worked}

\section{Practice Problems}

\begin{problems}
  \item Decide whether $S=\{(x,y)\in\R^{2}: x+2y\le 8\}$ is convex. Justify
        algebraically.
  \item Is the unit circle $S=\{(x,y): x^{2}+y^{2}=1\}$ (the boundary only,
        not the disk) convex?
  \item Show that the closed disk $D=\{(x,y): x^{2}+y^{2}\le 9\}$ is convex.
  \item Decide the convexity of $\{(x,y): y\ge x^{2}\}$.
  \item Decide the convexity of $\{(x,y): y\le -x^{2}\}$.
  \item Is the annulus $A=\{(x,y): 1\le x^{2}+y^{2}\le 4\}$ (a ring) convex?
  \item Is the union $\{x\in\R: x\le -1\}\cup\{x\in\R: x\ge 1\}$ convex?
  \item Is the open half-plane $\{(x,y): 3x-2y>5\}$ convex?
  \item Consider the budget set
        $B=\{(x_{1},x_{2}): 2x_{1}+3x_{2}\le 12,\ x_{1}\ge 0,\ x_{2}\ge 0\}$.
        Explain why $B$ is convex and find its three corner points.
  \item Sketch the polyhedron
        $P=\{(x,y): x\ge 0,\ y\ge 0,\ x+2y\le 6,\ 3x+y\le 9\}$ and list its
        vertices.
  \item Find the convex combination of $(1,0)$ and $(0,1)$ corresponding to
        $\lambda=0.3$ (with $\lambda$ multiplying the first point).
  \item Express the midpoint of $(2,5)$ and $(8,1)$ as a convex combination
        and compute it.
  \item Compute the convex combination of $(0,0)$, $(4,0)$, $(0,8)$ with
        weights $\tfrac12,\tfrac14,\tfrac14$, and verify that the resulting
        point lies in the triangle with these three vertices.
  \item Determine the convex hull of $\{(0,0),(2,0),(0,2),(1,1)\}$.
  \item Show that the intersection of any (possibly infinite) collection of
        convex sets is convex.
  \item Show by example that the union of two convex sets need not be
        convex.
  \item Show that if $S\subseteq\R^{n}$ is convex and $\alpha\in\R$, then
        $\alpha S=\{\alpha x : x\in S\}$ is convex.
  \item Show that the Minkowski sum
        $S_{1}+S_{2}=\{x+y : x\in S_{1},\ y\in S_{2}\}$ of two convex sets
        is convex.
  \item Is $\{(x,y): xy\ge 1,\ x>0\}$ convex?
  \item Is $\{(x,y): xy\le 1,\ x>0,\ y>0\}$ convex?
\end{problems}

\section{Solutions}

\begin{solutions}
  \item \textbf{Convex.} Take $(x_{1},y_{1}),(x_{2},y_{2})\in S$ and
        $\lambda\in[0,1]$. By linearity,
        \[
          \bigl(\lambda x_{1}+(1-\lambda)x_{2}\bigr)
          +2\bigl(\lambda y_{1}+(1-\lambda)y_{2}\bigr)
          =\lambda(x_{1}+2y_{1})+(1-\lambda)(x_{2}+2y_{2})
          \le 8\lambda+8(1-\lambda)=8 .
        \]
        The combination satisfies the defining inequality, so it lies in
        $S$: a half-plane is convex.

  \item \textbf{Not convex.} The points $(1,0)$ and $(-1,0)$ lie on the
        circle ($1^{2}+0^{2}=1$), but their midpoint
        $\tfrac12(1,0)+\tfrac12(-1,0)=(0,0)$ satisfies $0^{2}+0^{2}=0\neq
        1$, so it is not on the circle. The segment leaves the set
        immediately.

  \item Write $u=(x_{1},y_{1})$, $v=(x_{2},y_{2})$ with $\|u\|\le 3$,
        $\|v\|\le 3$ (membership in $D$ means the norm is at most
        $\sqrt{9}=3$). By the triangle inequality and homogeneity of the
        norm, for $\lambda\in[0,1]$,
        \[
          \|\lambda u+(1-\lambda)v\|
          \le \lambda\|u\|+(1-\lambda)\|v\|
          \le 3\lambda+3(1-\lambda)=3 ,
        \]
        so the combination has squared norm at most $9$ and lies in $D$.
        \textbf{Convex.}

  \item \textbf{Convex.} Let $y_{1}\ge x_{1}^{2}$ and $y_{2}\ge x_{2}^{2}$,
        and take $\lambda\in[0,1]$. Using the square inequality
        $(\lambda x_{1}+(1-\lambda)x_{2})^{2}\le\lambda
        x_{1}^{2}+(1-\lambda)x_{2}^{2}$ (the difference is
        $\lambda(1-\lambda)(x_{1}-x_{2})^{2}\ge 0$):
        \[
          \bigl(\lambda x_{1}+(1-\lambda)x_{2}\bigr)^{2}
          \le \lambda x_{1}^{2}+(1-\lambda)x_{2}^{2}
          \le \lambda y_{1}+(1-\lambda)y_{2} .
        \]
        So the combination again satisfies $y\ge x^{2}$. This is the region
        on and above the parabola, the graph shortcut with
        $g(x)=x^{2}$, $g''(x)=2\ge 0$.

  \item \textbf{Convex.} Let $y_{1}\le -x_{1}^{2}$ and $y_{2}\le
        -x_{2}^{2}$. Then
        \[
          \lambda y_{1}+(1-\lambda)y_{2}
          \le -\bigl(\lambda x_{1}^{2}+(1-\lambda)x_{2}^{2}\bigr)
          \le -\bigl(\lambda x_{1}+(1-\lambda)x_{2}\bigr)^{2},
        \]
        where the last step is the square inequality multiplied by $-1$.
        The region on and below the concave function $g(x)=-x^{2}$
        ($g''(x)=-2\le 0$) is convex, the graph shortcut again.

  \item \textbf{Not convex.} The points $(-\tfrac32,0)$ and $(\tfrac32,0)$
        lie in $A$ since $(\tfrac32)^{2}+0^{2}=\tfrac94$ and $1\le
        \tfrac94\le 4$. Their midpoint is $(0,0)$ with
        $0^{2}+0^{2}=0<1$, which violates the inner constraint. The hole in
        the middle destroys convexity.

  \item \textbf{Not convex.} Both $-1$ and $1$ belong to the union, but
        their midpoint $\tfrac12(-1)+\tfrac12(1)=0$ satisfies neither
        $0\le -1$ nor $0\ge 1$. The gap between the two pieces breaks the
        segment.

  \item \textbf{Convex.} For $\lambda\in(0,1)$ both weights are strictly
        positive, so from $3x_{1}-2y_{1}>5$ and $3x_{2}-2y_{2}>5$,
        \[
          3\bigl(\lambda x_{1}+(1-\lambda)x_{2}\bigr)
          -2\bigl(\lambda y_{1}+(1-\lambda)y_{2}\bigr)
          =\lambda(3x_{1}-2y_{1})+(1-\lambda)(3x_{2}-2y_{2})
          >5\lambda+5(1-\lambda)=5 .
        \]
        For $\lambda\in\{0,1\}$ the combination is one of the two original
        points. Open half-planes are convex just like closed ones.

  \item $B$ is the intersection of the three half-planes
        $\{2x_{1}+3x_{2}\le 12\}$, $\{x_{1}\ge 0\}$, $\{x_{2}\ge 0\}$. Each
        is convex, and intersections of convex sets are convex (Problem
        15), so $B$ is convex. Corner points: the origin $(0,0)$; on the
        $x_{1}$-axis, $x_{2}=0$ gives $2x_{1}=12$, i.e.\ $(6,0)$; on the
        $x_{2}$-axis, $x_{1}=0$ gives $3x_{2}=12$, i.e.\ $(0,4)$.
        Check: $2\cdot 6+3\cdot 0=12$ and $2\cdot 0+3\cdot 4=12$. $B$ is
        the triangle with vertices $(0,0)$, $(6,0)$, $(0,4)$.

  \item $P$ is the intersection of four half-planes, hence convex.
        Vertices:
        \begin{itemize}[leftmargin=1.6em,itemsep=0.3ex]
          \item origin: $(0,0)$;
          \item on the $x$-axis ($y=0$): the binding constraint is
                $3x+y=9$, so $x=3$, giving $(3,0)$; check $x+2y=3\le 6$;
                the alternative $x+2y=6$ would give $(6,0)$, but
                $3\cdot 6+0=18>9$ is infeasible;
          \item intersection of $x+2y=6$ and $3x+y=9$: from the first,
                $x=6-2y$; substituting, $3(6-2y)+y=18-5y=9$, so
                $y=\tfrac95$ and $x=6-\tfrac{18}{5}=\tfrac{12}{5}$, giving
                $\bigl(\tfrac{12}{5},\tfrac95\bigr)$;
          \item on the $y$-axis ($x=0$): the binding constraint is
                $x+2y=6$, so $y=3$, giving $(0,3)$; check $3\cdot 0+3=3\le
                9$; the alternative $3x+y=9$ would give $(0,9)$, but
                $0+2\cdot 9=18>6$ is infeasible.
        \end{itemize}
        $P$ is the quadrilateral with vertices $(0,0)$, $(3,0)$,
        $\bigl(\tfrac{12}{5},\tfrac95\bigr)$, $(0,3)$.

  \item $0.3\,(1,0)+0.7\,(0,1)=(0.3,\,0.7)$. The weights $0.3$ and $0.7$
        are nonnegative and sum to $1$, so this is a valid convex
        combination, a point three tenths of the way from $(0,1)$ to
        $(1,0)$.

  \item The midpoint is the convex combination with
        $\lambda=\tfrac12$:
        \[
          \tfrac12(2,5)+\tfrac12(8,1)
          =\Bigl(\tfrac{2+8}{2},\,\tfrac{5+1}{2}\Bigr)=(5,3) .
        \]

  \item The combination is
        \[
          \tfrac12(0,0)+\tfrac14(4,0)+\tfrac14(0,8)
          =(0+1+0,\;0+0+2)=(1,2),
        \]
        with weights $\tfrac12+\tfrac14+\tfrac14=1$, all nonnegative. The
        triangle with vertices $(0,0)$, $(4,0)$, $(0,8)$ is described by
        $x\ge 0$, $y\ge 0$, and $2x+y\le 8$ (the line through $(4,0)$ and
        $(0,8)$). At $(1,2)$: $1\ge 0$, $2\ge 0$, and $2\cdot 1+2=4\le 8$.
        The point lies inside the triangle, as every convex combination of
        the vertices must.

  \item The point $(1,1)$ is itself a convex combination of the others:
        $(1,1)=\tfrac12(2,0)+\tfrac12(0,2)$, the midpoint of the segment
        from $(2,0)$ to $(0,2)$. It therefore adds nothing to the hull:
        \[
          \operatorname{conv}\{(0,0),(2,0),(0,2),(1,1)\}
          =\operatorname{conv}\{(0,0),(2,0),(0,2)\},
        \]
        the closed triangle with vertices $(0,0)$, $(2,0)$, $(0,2)$
        (with $(1,1)$ on its hypotenuse).

  \item Let $\{S_{i}\}_{i\in I}$ be any family of convex sets and let
        $x,y\in\bigcap_{i\in I}S_{i}$, $\lambda\in[0,1]$. For each
        $i\in I$ we have $x,y\in S_{i}$, so by convexity of $S_{i}$,
        $\lambda x+(1-\lambda)y\in S_{i}$. Since this holds for every $i$,
        the combination lies in $\bigcap_{i\in I}S_{i}$. Hence the
        intersection is convex. (Nothing in the argument requires $I$ to be
        finite.)

  \item Take $S_{1}=\{x\in\R: x\le 0\}$ and $S_{2}=\{x\in\R: x\ge 1\}$,
        both convex half-lines. Their union contains $0$ and $1$, but the
        midpoint $\tfrac12\cdot 0+\tfrac12\cdot 1=\tfrac12$ satisfies
        neither $\tfrac12\le 0$ nor $\tfrac12\ge 1$, so it is not in the
        union. The union is not convex.

  \item Two generic elements of $\alpha S$ are $\alpha u$ and $\alpha v$
        with $u,v\in S$. For $\lambda\in[0,1]$,
        \[
          \lambda(\alpha u)+(1-\lambda)(\alpha v)
          =\alpha\bigl(\lambda u+(1-\lambda)v\bigr),
        \]
        and $\lambda u+(1-\lambda)v\in S$ by convexity of $S$, so the
        combination is $\alpha$ times an element of $S$, i.e.\ an element
        of $\alpha S$. Hence $\alpha S$ is convex.

  \item Two generic elements of $S_{1}+S_{2}$ are $u_{1}+v_{1}$ and
        $u_{2}+v_{2}$ with $u_{i}\in S_{1}$, $v_{i}\in S_{2}$. For
        $\lambda\in[0,1]$, regroup:
        \[
          \lambda(u_{1}+v_{1})+(1-\lambda)(u_{2}+v_{2})
          =\underbrace{\bigl[\lambda u_{1}+(1-\lambda)u_{2}\bigr]}_{\in
          S_{1}}
          +\underbrace{\bigl[\lambda v_{1}+(1-\lambda)v_{2}\bigr]}_{\in
          S_{2}} ,
        \]
        using convexity of $S_{1}$ and $S_{2}$. The result is again a sum
        of an element of $S_{1}$ and an element of $S_{2}$, so
        $S_{1}+S_{2}$ is convex.

  \item \textbf{Convex.} On $x>0$ the condition $xy\ge 1$ is equivalent to
        $y\ge \dfrac{1}{x}$, so the set is the region on and above the
        graph of $g(x)=\dfrac{1}{x}$ on $x>0$. Since
        $g''(x)=\dfrac{2}{x^{3}}>0$ for $x>0$, $g$ is convex and the graph
        shortcut applies. Direct verification with $\mu=1-\lambda$: for
        $(x_{1},y_{1}),(x_{2},y_{2})$ in the set,
        \[
          (\lambda x_{1}+\mu x_{2})(\lambda y_{1}+\mu y_{2})
          =\lambda^{2}x_{1}y_{1}+\mu^{2}x_{2}y_{2}
          +\lambda\mu\,(x_{1}y_{2}+x_{2}y_{1})
          \ge \lambda^{2}+\mu^{2}+2\lambda\mu=(\lambda+\mu)^{2}=1 ,
        \]
        because $x_{1}y_{1}\ge 1$, $x_{2}y_{2}\ge 1$, and by the
        AM--GM inequality
        $x_{1}y_{2}+x_{2}y_{1}\ge
        2\sqrt{(x_{1}y_{2})(x_{2}y_{1})}
        =2\sqrt{(x_{1}y_{1})(x_{2}y_{2})}\ge 2$. Also $\lambda
        x_{1}+\mu x_{2}>0$, so the combination stays in the set.

  \item \textbf{Not convex.} Take $(2,\,0.4)$ and $(0.4,\,2)$: both have
        product $2\cdot 0.4=0.8\le 1$ with positive coordinates, so both
        lie in the set. Their midpoint is $(1.2,\,1.2)$ with product
        $1.2\cdot 1.2=1.44>1$, which violates the constraint. The region
        under the hyperbola in the positive quadrant is not convex.
\end{solutions}

\section{Additional Problems}

\begin{problems}
  \item Decide the convexity of $\{(x,y,z)\in\R^{3}: x+y+z=1\}$ (a plane in
        $\R^{3}$).
  \item Show that $\{(x,y): y\ge |x|\}$ is convex.
  \item Is $\{(x,y): \mathrm{e}^{x}\le y\}$ convex?
  \item Is $\{(x,y): \ln x\ge y,\ x>0\}$ convex?
  \item Prove: if $f:\R^{n}\to\R$ is convex, then the lower contour set
        $\{x: f(x)\le c\}$ is convex for every $c\in\R$.
  \item Use the previous problem to conclude that
        $\{(x,y): x^{2}+y^{2}\le 4\}$ is convex.
  \item Prove: if $f:\R^{n}\to\R$ is concave, then the upper contour set
        $\{x: f(x)\ge c\}$ is convex for every $c\in\R$. Which weaker
        property of $f$ is defined by exactly this conclusion?
  \item Let $u(x_{1},x_{2})=x_{1}x_{2}$ on the positive quadrant
        $x_{1},x_{2}>0$. Show that every upper contour set
        $\{u\ge c\}$ with $c>0$ is convex, and conclude that $u$ is
        quasiconcave. Is $u$ concave?
  \item Give an example of a non-convex set in $\R^{2}$ whose convex hull
        is the entire plane.
  \item Show that the diamond $\{(x,y): |x|+|y|\le 1\}$ is convex by
        writing it as an intersection of four half-planes.

  \item Classify each plane figure as \emph{strictly convex},
        \emph{convex but not strictly convex}, or \emph{not convex}:
        (a) the closed disc $x^{2}+y^{2}\le4$;
        (b) the square $[0,1]\times[0,1]$;
        (c) the filled ellipse $\tfrac{x^{2}}{9}+\tfrac{y^{2}}{4}\le1$;
        (d) the half-disc $\{x^{2}+y^{2}\le1,\ y\ge0\}$;
        (e) the notched disc
        $\{x^{2}+y^{2}\le4\}\setminus\{(x,y): x>1,\ |y|<\tfrac12\}$;
        (f) the region $\{(x,y): y\ge x^{2}\}$.
  \item Classify the solids in $\R^{3}$ in the same way:
        (a) the ball $x^{2}+y^{2}+z^{2}\le1$;
        (b) the cube $[0,1]^{3}$;
        (c) the cylinder $\{x^{2}+y^{2}\le1,\ 0\le z\le2\}$;
        (d) the solid cone $\{\sqrt{x^{2}+y^{2}}\le z\le1\}$;
        (e) a solid torus (a ring doughnut).
  \item Prove that the closed disc $D=\{p\in\R^{2}: \|p\|\le r\}$ is
        strictly convex. (Hint: expand
        $\|\lambda p+(1-\lambda)q\|^{2}$.)
  \item Show that $U=\{(x,y): x>0,\ xy\ge4\}$ is strictly convex, whereas
        $V=\{(x,y): x\ge0,\ y\ge0,\ x+y\ge4\}$ is convex but not strictly
        convex. Which utility functions have these two sets as upper contour
        sets, and what does the difference mean for the uniqueness of the
        optimal bundle on a budget line?
\end{problems}

\section{Additional Solutions}

\begin{solutions}
  \item \textbf{Convex.} For two points of the plane and $\lambda\in[0,1]$,
        \[
          \bigl(\lambda x_{1}+(1-\lambda)x_{2}\bigr)
          +\bigl(\lambda y_{1}+(1-\lambda)y_{2}\bigr)
          +\bigl(\lambda z_{1}+(1-\lambda)z_{2}\bigr)
          =\lambda(x_{1}+y_{1}+z_{1})+(1-\lambda)(x_{2}+y_{2}+z_{2})
        \]
        \[
          =\lambda\cdot 1+(1-\lambda)\cdot 1=1 .
        \]
        The combination satisfies the equation, so the plane (an affine
        subspace) is convex.

  \item Let $y_{1}\ge |x_{1}|$ and $y_{2}\ge |x_{2}|$, and take
        $\lambda\in[0,1]$. By the triangle inequality for the absolute
        value,
        \[
          \bigl|\lambda x_{1}+(1-\lambda)x_{2}\bigr|
          \le \lambda|x_{1}|+(1-\lambda)|x_{2}|
          \le \lambda y_{1}+(1-\lambda)y_{2} ,
        \]
        so the combination again satisfies $y\ge|x|$. \textbf{Convex},
        the region above the V-shaped graph of the convex function $|x|$.

  \item \textbf{Convex.} The set is $\{y\ge \mathrm{e}^{x}\}$, the region
        on and above the graph of $g(x)=\mathrm{e}^{x}$, which is convex
        since $g''(x)=\mathrm{e}^{x}>0$. Directly: with $y_{1}\ge
        \mathrm{e}^{x_{1}}$, $y_{2}\ge \mathrm{e}^{x_{2}}$ and the
        convexity of the exponential,
        \[
          \lambda y_{1}+(1-\lambda)y_{2}
          \ge \lambda \mathrm{e}^{x_{1}}+(1-\lambda)\mathrm{e}^{x_{2}}
          \ge \mathrm{e}^{\lambda x_{1}+(1-\lambda)x_{2}} .
        \]

  \item \textbf{Convex.} The set is $\{y\le \ln x,\ x>0\}$, the region on
        and below the graph of $g(x)=\ln x$, which is concave since
        $g''(x)=-\dfrac{1}{x^{2}}<0$. Directly: $x>0$ is preserved by
        convex combination, and by concavity of $\ln$,
        \[
          \lambda y_{1}+(1-\lambda)y_{2}
          \le \lambda \ln x_{1}+(1-\lambda)\ln x_{2}
          \le \ln\bigl(\lambda x_{1}+(1-\lambda)x_{2}\bigr) .
        \]

  \item Let $f(x)\le c$, $f(y)\le c$, and $\lambda\in[0,1]$. By convexity
        of $f$,
        \[
          f\bigl(\lambda x+(1-\lambda)y\bigr)
          \le \lambda f(x)+(1-\lambda)f(y)
          \le \lambda c+(1-\lambda)c=c ,
        \]
        so the combination lies in $\{x: f(x)\le c\}$, which is therefore
        convex.

  \item The function $f(x,y)=x^{2}+y^{2}$ is convex as the sum of the
        convex functions $x^{2}$ and $y^{2}$ (each has second derivative
        $2>0$ in its own variable, and a sum of convex functions is
        convex). The set in question is the lower contour set
        $\{f\le 4\}$, so by the previous problem it is convex, the
        closed disk of radius $2$.

  \item Let $f(x)\ge c$, $f(y)\ge c$, and $\lambda\in[0,1]$. By concavity
        of $f$,
        \[
          f\bigl(\lambda x+(1-\lambda)y\bigr)
          \ge \lambda f(x)+(1-\lambda)f(y)
          \ge \lambda c+(1-\lambda)c=c ,
        \]
        so the upper contour set is convex. The property ``every upper
        contour set of $f$ is convex'' is precisely
        \textbf{quasiconcavity}; the computation shows that every concave
        function is quasiconcave (the converse fails, as the next problem
        shows).

  \item For $c>0$ and $x_{1}>0$, the condition $x_{1}x_{2}\ge c$ is
        equivalent to $x_{2}\ge \dfrac{c}{x_{1}}$, so
        \[
          \{u\ge c\}=\Bigl\{(x_{1},x_{2}): x_{2}\ge
          \tfrac{c}{x_{1}},\ x_{1}>0\Bigr\},
        \]
        the region on and above the graph of
        $g(x_{1})=\dfrac{c}{x_{1}}$ on $x_{1}>0$. Since
        $g''(x_{1})=\dfrac{2c}{x_{1}^{3}}>0$ for $c>0$, $g$ is convex and
        the set is convex (same argument as Practice Problem~19, which is
        the case $c=1$). For $c\le 0$ the upper contour set is the whole
        positive quadrant, also convex. Every upper contour set is convex,
        so $u$ is \textbf{quasiconcave}. But $u$ is \emph{not} concave:
        along the line $x_{1}=x_{2}=t$ ($t>0$) we get $u=t^{2}$, a strictly
        convex function of $t$, whereas a concave function must remain
        concave along every line. Quasiconcavity is strictly weaker than
        concavity.

  \item Take $S=\R^{2}\setminus\{(0,0)\}$, the plane with the origin
        removed. It is not convex: $(-1,0)$ and $(1,0)$ lie in $S$, but
        their midpoint $(0,0)$ does not. Its convex hull contains every
        point of $S$ and also the origin, since
        $(0,0)=\tfrac12(-1,0)+\tfrac12(1,0)$ is a convex combination of
        points of $S$. Hence $\operatorname{conv}(S)=\R^{2}$.
        (Another example: $\mathbb{Q}^{2}$, the points with rational
        coordinates, is not convex, yet its convex hull is all of
        $\R^{2}$.)

  \item Since $|x|=\max(x,-x)$ and $|y|=\max(y,-y)$, the condition
        $|x|+|y|\le 1$ holds if and only if all four sign combinations
        satisfy it:
        \[
          x+y\le 1,\qquad x-y\le 1,\qquad -x+y\le 1,\qquad -x-y\le 1 .
        \]
        Each of these describes a half-plane, which is convex, and the
        diamond is the intersection of the four. Intersections of convex
        sets are convex, so the diamond is convex, the square with
        vertices $(1,0)$, $(0,1)$, $(-1,0)$, $(0,-1)$.

  \item (a) \textbf{Strictly convex}: the boundary circle contains no
        segment (proof in Additional Problem~13).
        (b) \textbf{Convex, not strictly}: $(0,0)$ and $(1,0)$ lie in the
        square, but their midpoint $(\tfrac12,0)$ lies on the bottom edge,
        not in the interior.
        (c) \textbf{Strictly convex}: the ellipse is the image of a disc
        under the stretch $(u,v)\mapsto(3u,2v)$, which maps chords to chords
        and interior points to interior points.
        (d) \textbf{Convex, not strictly}: $(-1,0)$ and $(1,0)$ belong to it
        and their midpoint $(0,0)$ sits on the flat diameter, a boundary
        point.
        (e) \textbf{Not convex}: $(1.5,\,0.6)$ and $(1.5,\,-0.6)$ satisfy
        $2.25+0.36=2.61\le4$ and $|y|=0.6\ge\tfrac12$, so both belong to the
        set, but their midpoint $(1.5,0)$ lies in the removed notch.
        (f) \textbf{Strictly convex}: take two different points
        $(a,s)$, $(b,t)$ with $s\ge a^{2}$, $t\ge b^{2}$ and
        $\lambda\in(0,1)$. If $a\neq b$, the square inequality is strict,
        $\bigl(\lambda a+(1-\lambda)b\bigr)^{2}<\lambda a^{2}+(1-\lambda)b^{2}
        \le\lambda s+(1-\lambda)t$; if $a=b$ then $s\neq t$, so
        $\lambda s+(1-\lambda)t>\min\{s,t\}\ge a^{2}$. Either way the
        combination satisfies $y>x^{2}$ strictly, which describes an open
        region, hence an interior point.
  \item (a) \textbf{Strictly convex}: the same norm argument as for the
        disc (Additional Problem~13) works in any dimension.
        (b) \textbf{Convex, not strictly}: the midpoint of $(0,0,0)$ and
        $(1,0,0)$ lies on an edge.
        (c) \textbf{Convex, not strictly}: it is the intersection of the
        convex set $\{x^{2}+y^{2}\le1\}$ with two half-spaces, but
        $(1,0,0)$ and $(1,0,2)$ lie on the side wall and so does their
        midpoint $(1,0,1)$; the top and bottom discs are flat as well.
        (d) \textbf{Convex, not strictly}: $\sqrt{x^{2}+y^{2}}$ is a convex
        function (a norm), so $\{\sqrt{x^{2}+y^{2}}\le z\}$ is convex, and
        intersecting with $z\le1$ keeps convexity. But the segment from the
        tip $(0,0,0)$ to $(1,0,1)$ lies on the mantle: its midpoint
        $(\tfrac12,0,\tfrac12)$ satisfies $\sqrt{x^{2}+y^{2}}=z$, a boundary
        point.
        (e) \textbf{Not convex}: two points on opposite sides of the ring
        have their midpoint in the hole.
  \item Let $p\neq q$ with $\|p\|\le r$, $\|q\|\le r$ and
        $\lambda\in(0,1)$. Expanding the inner product,
        \[
          \|\lambda p+(1-\lambda)q\|^{2}
          =\lambda\|p\|^{2}+(1-\lambda)\|q\|^{2}
           -\lambda(1-\lambda)\|p-q\|^{2},
        \]
        which is the vector form of the square inequality. The first two
        terms are at most $\lambda r^{2}+(1-\lambda)r^{2}=r^{2}$, and the
        last term is strictly positive because $p\neq q$. Hence
        $\|\lambda p+(1-\lambda)q\|<r$: the combination lies in the open
        disc, i.e.\ in the interior of $D$. So $D$ is \textbf{strictly
        convex}.
  \item \emph{$U$.} Write $U=\{x>0,\ y\ge g(x)\}$ with $g(x)=4/x$,
        $g''(x)=8/x^{3}>0$: $g$ is strictly convex. Take
        $P=(x_1,y_1)\neq Q=(x_2,y_2)$ in $U$ and $\lambda\in(0,1)$. Then
        \[
          \lambda y_1+(1-\lambda)y_2
          \;\ge\;\lambda g(x_1)+(1-\lambda)g(x_2)
          \;\ge\; g\bigl(\lambda x_1+(1-\lambda)x_2\bigr).
        \]
        If $x_1\neq x_2$ the second inequality is strict ($g$ strictly
        convex); if $x_1=x_2$ then $y_1\neq y_2$ and the first is strict
        unless $y_1=y_2=g(x_1)$, which is excluded. So the combination
        satisfies $xy>4$ strictly, an open condition: it is an interior
        point and $U$ is \textbf{strictly convex}.

        \emph{$V$.} It is an intersection of three half-planes, so convex.
        But $(4,0)$ and $(0,4)$ lie on the boundary line $x+y=4$ and so does
        their midpoint $(2,2)$: \textbf{not strictly convex}.

        $U$ is the upper contour set $\{u\ge4\}$ of $u(x,y)=xy$ (on $x>0$);
        $V$ is $\{u\ge4\}$ of the perfect-substitutes utility $u=x+y$ on
        the quadrant. On a budget line the optimum touches the highest
        reachable upper contour set. A strictly convex set can touch a line
        in only one point, so $xy$ has a unique optimal bundle; the flat
        edge of $V$ can lie \emph{along} a budget line with $p_x=p_y$, and
        then every bundle on the line is optimal.
\end{solutions}
